\IfFileExists{stacks-project.cls}{%
\documentclass{stacks-project}
}{%
\documentclass{amsart}
}

% For dealing with references we use the comment environment
\usepackage{verbatim}
\newenvironment{reference}{\comment}{\endcomment}
%\newenvironment{reference}{}{}
\newenvironment{slogan}{\comment}{\endcomment}
\newenvironment{history}{\comment}{\endcomment}

% For commutative diagrams we use Xy-pic
\usepackage[all]{xy}

% We use 2cell for 2-commutative diagrams.
\xyoption{2cell}
\UseAllTwocells

% We use multicol for the list of chapters between chapters
\usepackage{multicol}

% This is generally recommended for better output
\usepackage{lmodern}
\usepackage[T1]{fontenc}

% For cross-file-references
\usepackage{xr-hyper}

% Package for hypertext links:
\usepackage{hyperref}

% For any local file, say "hello.tex" you want to link to please
% use \externaldocument[hello-]{hello}
\externaldocument[introduction-]{introduction}
\externaldocument[conventions-]{conventions}
\externaldocument[sets-]{sets}
\externaldocument[categories-]{categories}
\externaldocument[topology-]{topology}
\externaldocument[sheaves-]{sheaves}
\externaldocument[sites-]{sites}
\externaldocument[stacks-]{stacks}
\externaldocument[fields-]{fields}
\externaldocument[algebra-]{algebra}
\externaldocument[brauer-]{brauer}
\externaldocument[homology-]{homology}
\externaldocument[derived-]{derived}
\externaldocument[simplicial-]{simplicial}
\externaldocument[more-algebra-]{more-algebra}
\externaldocument[smoothing-]{smoothing}
\externaldocument[modules-]{modules}
\externaldocument[sites-modules-]{sites-modules}
\externaldocument[injectives-]{injectives}
\externaldocument[cohomology-]{cohomology}
\externaldocument[sites-cohomology-]{sites-cohomology}
\externaldocument[dga-]{dga}
\externaldocument[dpa-]{dpa}
\externaldocument[sdga-]{sdga}
\externaldocument[hypercovering-]{hypercovering}
\externaldocument[schemes-]{schemes}
\externaldocument[constructions-]{constructions}
\externaldocument[properties-]{properties}
\externaldocument[morphisms-]{morphisms}
\externaldocument[coherent-]{coherent}
\externaldocument[divisors-]{divisors}
\externaldocument[limits-]{limits}
\externaldocument[varieties-]{varieties}
\externaldocument[topologies-]{topologies}
\externaldocument[descent-]{descent}
\externaldocument[perfect-]{perfect}
\externaldocument[more-morphisms-]{more-morphisms}
\externaldocument[flat-]{flat}
\externaldocument[groupoids-]{groupoids}
\externaldocument[more-groupoids-]{more-groupoids}
\externaldocument[etale-]{etale}
\externaldocument[chow-]{chow}
\externaldocument[intersection-]{intersection}
\externaldocument[pic-]{pic}
\externaldocument[weil-]{weil}
\externaldocument[adequate-]{adequate}
\externaldocument[dualizing-]{dualizing}
\externaldocument[duality-]{duality}
\externaldocument[discriminant-]{discriminant}
\externaldocument[derham-]{derham}
\externaldocument[local-cohomology-]{local-cohomology}
\externaldocument[algebraization-]{algebraization}
\externaldocument[curves-]{curves}
\externaldocument[resolve-]{resolve}
\externaldocument[models-]{models}
\externaldocument[functors-]{functors}
\externaldocument[equiv-]{equiv}
\externaldocument[pione-]{pione}
\externaldocument[etale-cohomology-]{etale-cohomology}
\externaldocument[proetale-]{proetale}
\externaldocument[relative-cycles-]{relative-cycles}
\externaldocument[more-etale-]{more-etale}
\externaldocument[trace-]{trace}
\externaldocument[crystalline-]{crystalline}
\externaldocument[spaces-]{spaces}
\externaldocument[spaces-properties-]{spaces-properties}
\externaldocument[spaces-morphisms-]{spaces-morphisms}
\externaldocument[decent-spaces-]{decent-spaces}
\externaldocument[spaces-cohomology-]{spaces-cohomology}
\externaldocument[spaces-limits-]{spaces-limits}
\externaldocument[spaces-divisors-]{spaces-divisors}
\externaldocument[spaces-over-fields-]{spaces-over-fields}
\externaldocument[spaces-topologies-]{spaces-topologies}
\externaldocument[spaces-descent-]{spaces-descent}
\externaldocument[spaces-perfect-]{spaces-perfect}
\externaldocument[spaces-more-morphisms-]{spaces-more-morphisms}
\externaldocument[spaces-flat-]{spaces-flat}
\externaldocument[spaces-groupoids-]{spaces-groupoids}
\externaldocument[spaces-more-groupoids-]{spaces-more-groupoids}
\externaldocument[bootstrap-]{bootstrap}
\externaldocument[spaces-pushouts-]{spaces-pushouts}
\externaldocument[spaces-chow-]{spaces-chow}
\externaldocument[groupoids-quotients-]{groupoids-quotients}
\externaldocument[spaces-more-cohomology-]{spaces-more-cohomology}
\externaldocument[spaces-simplicial-]{spaces-simplicial}
\externaldocument[spaces-duality-]{spaces-duality}
\externaldocument[formal-spaces-]{formal-spaces}
\externaldocument[restricted-]{restricted}
\externaldocument[spaces-resolve-]{spaces-resolve}
\externaldocument[formal-defos-]{formal-defos}
\externaldocument[defos-]{defos}
\externaldocument[cotangent-]{cotangent}
\externaldocument[examples-defos-]{examples-defos}
\externaldocument[algebraic-]{algebraic}
\externaldocument[examples-stacks-]{examples-stacks}
\externaldocument[stacks-sheaves-]{stacks-sheaves}
\externaldocument[criteria-]{criteria}
\externaldocument[artin-]{artin}
\externaldocument[quot-]{quot}
\externaldocument[stacks-properties-]{stacks-properties}
\externaldocument[stacks-morphisms-]{stacks-morphisms}
\externaldocument[stacks-limits-]{stacks-limits}
\externaldocument[stacks-cohomology-]{stacks-cohomology}
\externaldocument[stacks-perfect-]{stacks-perfect}
\externaldocument[stacks-introduction-]{stacks-introduction}
\externaldocument[stacks-more-morphisms-]{stacks-more-morphisms}
\externaldocument[stacks-geometry-]{stacks-geometry}
\externaldocument[moduli-]{moduli}
\externaldocument[moduli-curves-]{moduli-curves}
\externaldocument[examples-]{examples}
\externaldocument[exercises-]{exercises}
\externaldocument[guide-]{guide}
\externaldocument[desirables-]{desirables}
\externaldocument[coding-]{coding}
\externaldocument[obsolete-]{obsolete}
\externaldocument[fdl-]{fdl}
\externaldocument[index-]{index}

% Theorem environments.
%
\theoremstyle{plain}
\newtheorem{theorem}[subsection]{Theorem}
\newtheorem{proposition}[subsection]{Proposition}
\newtheorem{lemma}[subsection]{Lemma}

\theoremstyle{definition}
\newtheorem{definition}[subsection]{Definition}
\newtheorem{example}[subsection]{Example}
\newtheorem{exercise}[subsection]{Exercise}
\newtheorem{situation}[subsection]{Situation}

\theoremstyle{remark}
\newtheorem{remark}[subsection]{Remark}
\newtheorem{remarks}[subsection]{Remarks}

\numberwithin{equation}{subsection}

% Macros
%
\def\lim{\mathop{\mathrm{lim}}\nolimits}
\def\colim{\mathop{\mathrm{colim}}\nolimits}
\def\Spec{\mathop{\mathrm{Spec}}}
\def\Hom{\mathop{\mathrm{Hom}}\nolimits}
\def\Ext{\mathop{\mathrm{Ext}}\nolimits}
\def\SheafHom{\mathop{\mathcal{H}\!\mathit{om}}\nolimits}
\def\SheafExt{\mathop{\mathcal{E}\!\mathit{xt}}\nolimits}
\def\Sch{\mathit{Sch}}
\def\Mor{\mathop{\mathrm{Mor}}\nolimits}
\def\Ob{\mathop{\mathrm{Ob}}\nolimits}
\def\Sh{\mathop{\mathit{Sh}}\nolimits}
\def\NL{\mathop{N\!L}\nolimits}
\def\CH{\mathop{\mathrm{CH}}\nolimits}
\def\proetale{{pro\text{-}\acute{e}tale}}
\def\etale{{\acute{e}tale}}
\def\QCoh{\mathit{QCoh}}
\def\Ker{\mathop{\mathrm{Ker}}}
\def\Im{\mathop{\mathrm{Im}}}
\def\Coker{\mathop{\mathrm{Coker}}}
\def\Coim{\mathop{\mathrm{Coim}}}

% Boxtimes
%
\DeclareMathSymbol{\boxtimes}{\mathbin}{AMSa}{"02}

%
% Macros for moduli stacks/spaces
%
\def\QCohstack{\mathcal{QC}\!\mathit{oh}}
\def\Cohstack{\mathcal{C}\!\mathit{oh}}
\def\Spacesstack{\mathcal{S}\!\mathit{paces}}
\def\Quotfunctor{\mathrm{Quot}}
\def\Hilbfunctor{\mathrm{Hilb}}
\def\Curvesstack{\mathcal{C}\!\mathit{urves}}
\def\Polarizedstack{\mathcal{P}\!\mathit{olarized}}
\def\Complexesstack{\mathcal{C}\!\mathit{omplexes}}
% \Pic is the operator that assigns to X its picard group, usage \Pic(X)
% \Picardstack_{X/B} denotes the Picard stack of X over B
% \Picardfunctor_{X/B} denotes the Picard functor of X over B
\def\Pic{\mathop{\mathrm{Pic}}\nolimits}
\def\Picardstack{\mathcal{P}\!\mathit{ic}}
\def\Picardfunctor{\mathrm{Pic}}
\def\Deformationcategory{\mathcal{D}\!\mathit{ef}}
\usepackage{cleveref}
\crefname{originallemma}{original lemma}{original lemmas}
\Crefname{originallemma}{Original lemma}{Original lemmas}
\makeatletter
\newlabel{derived-lemma-unbounded-right-derived}{{07K7}{source}{}{sourcepointer-1}{}}
\newlabel{derived-lemma-unbounded-right-derived@cref}{{[originallemma][0][]07K7}{[1][0][]source}{}{}{}}
\makeatother
\crefname{originallemma}{original lemma}{original lemmas}
\Crefname{originallemma}{Original lemma}{Original lemmas}
\makeatletter
\newlabel{etale-cohomology-lemma-cohomological-dimension-proper}{{095U}{source}{}{sourcepointer-2}{}}
\newlabel{etale-cohomology-lemma-cohomological-dimension-proper@cref}{{[originallemma][0][]095U}{[1][0][]source}{}{}{}}
\makeatother
\crefname{originallemma}{original lemma}{original lemmas}
\Crefname{originallemma}{Original lemma}{Original lemmas}
\makeatletter
\newlabel{etale-cohomology-lemma-proper-base-change-stalk}{{0DDF}{source}{}{sourcepointer-3}{}}
\newlabel{etale-cohomology-lemma-proper-base-change-stalk@cref}{{[originallemma][0][]0DDF}{[1][0][]source}{}{}{}}
\makeatother
\crefname{originallemma}{original lemma}{original lemmas}
\Crefname{originallemma}{Original lemma}{Original lemmas}
\makeatletter
\newlabel{etale-cohomology-lemma-supported-in-closed-points}{{0F1F}{source}{}{sourcepointer-4}{}}
\newlabel{etale-cohomology-lemma-supported-in-closed-points@cref}{{[originallemma][0][]0F1F}{[1][0][]source}{}{}{}}
\makeatother
\crefname{originallemma}{original lemma}{original lemmas}
\Crefname{originallemma}{Original lemma}{Original lemmas}
\makeatletter
\newlabel{lemma-lqf-f-shriek-separated-colimits}{{0F57}{source}{}{sourcepointer-5}{}}
\newlabel{lemma-lqf-f-shriek-separated-colimits@cref}{{[originallemma][0][]0F57}{[1][0][]source}{}{}{}}
\makeatother
\crefname{originallemma}{original lemma}{original lemmas}
\Crefname{originallemma}{Original lemma}{Original lemmas}
\makeatletter
\newlabel{lemma-lqf-f-shriek-stalk}{{0F5F}{source}{}{sourcepointer-6}{}}
\newlabel{lemma-lqf-f-shriek-stalk@cref}{{[originallemma][0][]0F5F}{[1][0][]source}{}{}{}}
\makeatother
\crefname{originallemma}{original lemma}{original lemmas}
\Crefname{originallemma}{Original lemma}{Original lemmas}
\makeatletter
\newlabel{lemma-lqf-separated-shriek-composition}{{0F79}{source}{}{sourcepointer-7}{}}
\newlabel{lemma-lqf-separated-shriek-composition@cref}{{[originallemma][0][]0F79}{[1][0][]source}{}{}{}}
\makeatother
\crefname{originallemma}{original lemma}{original lemmas}
\Crefname{originallemma}{Original lemma}{Original lemmas}
\makeatletter
\newlabel{lemma-proper-f-shriek}{{0F51}{source}{}{sourcepointer-8}{}}
\newlabel{lemma-proper-f-shriek@cref}{{[originallemma][0][]0F51}{[1][0][]source}{}{}{}}
\makeatother
\crefname{originallemma}{original lemma}{original lemmas}
\Crefname{originallemma}{Original lemma}{Original lemmas}
\makeatletter
\newlabel{lemma-shriek-proper-and-open}{{0F7B}{source}{}{sourcepointer-9}{}}
\newlabel{lemma-shriek-proper-and-open@cref}{{[originallemma][0][]0F7B}{[1][0][]source}{}{}{}}
\makeatother
\crefname{originallemma}{original lemma}{original lemmas}
\Crefname{originallemma}{Original lemma}{Original lemmas}
\makeatletter
\newlabel{morphisms-lemma-morphism-finite-type-bounded-dimension}{{0A3V}{source}{}{sourcepointer-10}{}}
\newlabel{morphisms-lemma-morphism-finite-type-bounded-dimension@cref}{{[originallemma][0][]0A3V}{[1][0][]source}{}{}{}}
\makeatother
\title{Canonical map from lower shriek along a locally quasi-finite composite to proper pushforward of lower shriek, and its invertibility}
\author{The Stacks Project Authors}
\date{stacks-project:more-etale.tex:local-snapshot-20261012; GNU FDL 1.2 or later}
\begin{document}
\maketitle
\setcounter{subsection}{1}
\noindent\textit{Editorial note.} All delivered mathematical text is the author's own, copied from the pinned Stacks Project chapter file more-etale.tex; no statement or proof step was rewritten, reordered or removed. The counted span is the author environment and its own proof environment. References inside the span are handed to the assembler as navigation stubs that carry the canonical tag of the cited tag only; whether each of them is load-bearing is recorded on the stub itself.\par\medskip
\noindent\textit{Editorial note.} Disclosure: two soft spots. (1) The author notes the separated case is a special case of Lemma lemma-shriek-proper-and-open and adds 'We urge the reader to skip the proof in the general case' (proof lines 3198-3201); the general-case proof is nonetheless given in full. (2) Case (a) is not re-proven: 'Case (a) follows from case (b) by the exact same argument as used in the proof of Lemma lemma-shriek-proper-and-open (using Lemma lemma-lqf-f-shriek-stalk instead of Lemma lemma-lqf-f-shriek-separated-colimits)' (proof lines 3251-3254) -- a parallel-case pointer to another proof with the substituted lemma named. The construction itself is fully displayed: the first equality uses lemma-lqf-separated-shriek-composition, whose hypotheses (upper morphism separated and locally of finite type) are met by the proper h. No verbatim 'we omit' occurs in the proof.\par\medskip
\noindent\textit{Allowed standard prerequisites (editorial).}\begin{itemize}
\item Composition isomorphism for lower shriek (Lemma lemma-lqf-separated-shriek-composition, canonical tag 0F79). cited at proof lines 3216-3218. Gives the first equality g\_! = h\_! f\_!: its hypotheses (upper morphism separated and locally of finite type, composite locally quasi-finite) are met since h is proper.
\item Proper lower shriek identification (Lemma lemma-proper-f-shriek, canonical tag 0F51). cited at proof lines 3216-3218. Gives the second equality h\_! = h\_* for the proper morphism h.
\item Vanishing for sheaves supported in closed points (Etale Cohomology, Lemma etale-cohomology-lemma-supported-in-closed-points, canonical tag 0F1F). cited at proof lines 3244-3248. The terminal vanishing: cohomology of a sheaf supported on the closed points of the proper fibre vanishes in positive degrees.
\end{itemize}\medskip
\ifdefined\hypertarget\hypertarget{sourcepointer-1}{}\fi
\noindent\textit{Source reference.} Original lemma 07K7: sources/stacks-chow/tags.txt, source lines 9647--9647. Cited at proof lines 3230-3233: with finite cohomological dimension, vanishing of the higher derived pushforwards makes the map of complexes a quasi-isomorphism.\par\medskip
\ifdefined\hypertarget\hypertarget{sourcepointer-2}{}\fi
\noindent\textit{Source reference.} Original lemma 095U: sources/stacks-chow/tags.txt, source lines 11597--11597. Cited at proof lines 3227-3229: Rh\_* has finite cohomological dimension once the fibre dimension is bounded.\par\medskip
\ifdefined\hypertarget\hypertarget{sourcepointer-3}{}\fi
\noindent\textit{Source reference.} Original lemma 0DDF: sources/stacks-chow/tags.txt, source lines 16765--16765. Cited at proof lines 3242-3243: reduces the vanishing to the cohomology of the geometric fibre.\par\medskip
\ifdefined\hypertarget\hypertarget{sourcepointer-4}{}\fi
\noindent\textit{Source reference.} Original lemma 0F1F: sources/stacks-chow/tags.txt, source lines 18795--18795. Cited at proof lines 3247-3248: the required vanishing for sheaves supported on the closed points of the proper fibre.\par\medskip
\ifdefined\hypertarget\hypertarget{sourcepointer-5}{}\fi
\noindent\textit{Source reference.} Original lemma 0F57: sources/stacks-chow/tags.txt, source lines 18927--18927. Cited at proof line 3254: the lemma replaced by lemma-lqf-f-shriek-stalk when rerunning the argument of the earlier exchange lemma in case (a).\par\medskip
\ifdefined\hypertarget\hypertarget{sourcepointer-6}{}\fi
\noindent\textit{Source reference.} Original lemma 0F5F: sources/stacks-chow/tags.txt, source lines 18935--18935. Cited at proof lines 3239-3240 and 3253: stalks of f\_!F are nonzero only over points with finite residue field extension, the input to the vanishing argument.\par\medskip
\ifdefined\hypertarget\hypertarget{sourcepointer-7}{}\fi
\noindent\textit{Source reference.} Original lemma 0F79: sources/stacks-chow/tags.txt, source lines 18999--18999. Cited at proof lines 3216-3218: the first equality g\_! = h\_!f\_! of the construction.\par\medskip
\ifdefined\hypertarget\hypertarget{sourcepointer-8}{}\fi
\noindent\textit{Source reference.} Original lemma 0F51: sources/stacks-chow/tags.txt, source lines 18921--18921. Cited at proof lines 3216-3218: the second equality h\_! = h\_* of the construction.\par\medskip
\ifdefined\hypertarget\hypertarget{sourcepointer-9}{}\fi
\noindent\textit{Source reference.} Original lemma 0F7B: sources/stacks-chow/tags.txt, source lines 19001--19001. Cited at proof lines 3198-3199 and 3251-3252: the separated case is its special case, and case (a) reruns its argument.\par\medskip
\ifdefined\hypertarget\hypertarget{sourcepointer-10}{}\fi
\noindent\textit{Source reference.} Original lemma 0A3V: sources/stacks-chow/tags.txt, source lines 12753--12753. Cited at proof lines 3224-3226: working locally on Z bounds the dimension of the fibres of h.\par\medskip
\begin{lemma}
\label{lemma-shriek-lqf-and-proper}
Consider a commutative diagram of schemes
$$
\xymatrix{
X \ar[r]_f \ar[rd]_g & Y \ar[d]^h \\
& Z
}
$$
with $f$ and $g$ locally quasi-finite and $h$ proper. Let $\Lambda$ be a ring.
Funtorially in $K \in D(X_\etale, \Lambda)$ there is a canonical map
$$
g_!K \longrightarrow Rh_*(f_!K)
$$
in $D(Z_\etale, \Lambda)$. This map is an isomorphism if
(a) $K$ is bounded below and has torsion cohomology sheaves, or
(b) $\Lambda$ is a torsion ring.
\end{lemma}

\begin{proof}
This is a special case of Lemma \ref{lemma-shriek-proper-and-open}
if $f$ and $g$ are separated. We urge the reader to skip the proof
in the general case
as we'll mainly use the case where $f$ and $g$ are separated.

\medskip\noindent
Represent $K$ by a complex $\mathcal{K}^\bullet$ of sheaves of
$\Lambda$-modules on $X_\etale$. Choose a quasi-isomorphism
$f_!\mathcal{K}^\bullet \to \mathcal{I}^\bullet$ into a K-injective
complex $\mathcal{I}^\bullet$ of sheaves of $\Lambda$-modules on $Y_\etale$.
Consider the map
$$
g_!\mathcal{K}^\bullet =
h_!f_!\mathcal{K}^\bullet =
h_*f_!\mathcal{K}^\bullet
\longrightarrow
h_*\mathcal{I}^\bullet
$$
where the equalities are
Lemmas \ref{lemma-lqf-separated-shriek-composition}
and \ref{lemma-proper-f-shriek}. This map of complexes determines the map
$g_!K \to Rh_*(f_!K)$ of the statement of the lemma.

\medskip\noindent
Assume $\Lambda$ is torsion, i.e., we are in case (b). To check the map
is an isomorphism we may work locally on $Z$.
Hence we may assume that the dimension of fibres of $h$ is bounded,
see Morphisms, Lemma
\ref{morphisms-lemma-morphism-finite-type-bounded-dimension}.
Then we see that $Rh_*$ has finite cohomological dimension, see
\'Etale Cohomology, Lemma
\ref{etale-cohomology-lemma-cohomological-dimension-proper}.
Hence by Derived Categories, Lemma \ref{derived-lemma-unbounded-right-derived},
if we show that $R^qh_*(f_!\mathcal{F}) = 0$ for $q > 0$
and any sheaf $\mathcal{F}$ of $\Lambda$-modules on $X_\etale$, then
$h_*f_!\mathcal{K}^\bullet \to h_*\mathcal{I}^\bullet$ is a quasi-isomorphism.

\medskip\noindent
Observe that $\mathcal{G} = f_!\mathcal{F}$ is a sheaf of $\Lambda$-modules
on $Y$ whose stalks are nonzero only at points $y \in Y$ such that
$\kappa(y)/\kappa(h(y))$ is a finite extension. This follows from the
description of stalks of $f_!\mathcal{F}$ in
Lemma \ref{lemma-lqf-f-shriek-stalk}
and the fact that both $f$ and $g$ are locally quasi-finite.
Hence by the proper base change theorem (\'Etale Cohomology, Lemma
\ref{etale-cohomology-lemma-proper-base-change-stalk})
it suffices to show that $H^q(Y_{\overline{z}}, \mathcal{H}) = 0$
where $\mathcal{H}$ is a sheaf on the proper scheme $Y_{\overline{z}}$
over $\kappa(\overline{z})$ whose support is contained in the set
of closed points. Thus the required vanishing by \'Etale Cohomology, Lemma
\ref{etale-cohomology-lemma-supported-in-closed-points}.

\medskip\noindent
Case (a) follows from case (b) by the exact same argument as used
in the proof of Lemma \ref{lemma-shriek-proper-and-open}
(using Lemma \ref{lemma-lqf-f-shriek-stalk} instead of
Lemma \ref{lemma-lqf-f-shriek-separated-colimits}).
\end{proof}

\end{document}
